% Halaman judul penuh dan pernyataan hak cipta, wajib menurut
% Permendikbudristek 22/2022 Lampiran III Tabel 1. Bagian bertanda merah diisi
% setelah penerbit ditetapkan.
\cleardoublepage
\thispagestyle{empty}
\begin{center}
  \vspace*{0.12\textheight}
  {\HUGE \textcolor{praditagreen}{\textbf{Rekayasa\\[6pt] Perangkat Lunak\\[6pt] Berbasis Model\\[6pt] untuk DevOps}}}\\[0.06\textheight]
  {\Large Dari Sistem yang Berjalan sampai Generasi Artefak}\\[0.08\textheight]
  {\LARGE \textbf{Alfa Yohannis}}\\[0.22\textheight]
  {\Large \isian{nama penerbit}}\\[0.01\textheight]
  {\large \isian{tahun terbit}}
\end{center}
\clearpage

\thispagestyle{empty}
\vspace*{\fill}
\noindent\textbf{Rekayasa Perangkat Lunak Berbasis Model untuk DevOps}

\medskip
\noindent Buku ini disusun sebagai rujukan rekayasa perangkat lunak berbasis model dan DevOps, sekaligus sebagai pegangan mata kuliah IT30213 pada Program Studi Informatika, Universitas Pradita.

\medskip
\noindent
\begin{tabular}{@{}ll@{}}
  Penulis & Alfa Yohannis \\
  Editor & \isian{diisi penerbit} \\
  Desainer sampul & \isian{diisi penerbit} \\
  Penata letak & \isian{diisi penerbit} \\
  Ukuran & 17,6 $\times$ 25 cm (B5) \\
  Edisi & Cetakan pertama, \isian{tahun terbit} \\
  ISBN & \isian{diisi penerbit} \\
  Penerbit & \isian{nama dan alamat penerbit} \\
\end{tabular}

\medskip
\noindent Hak cipta \textcopyright{} \isian{tahun terbit} pada \isian{pemegang hak cipta, dikonfirmasi penerbit}.

\noindent Hak cipta dilindungi undang-undang. Dilarang memperbanyak sebagian atau seluruh isi buku ini dalam bentuk apa pun tanpa izin tertulis dari pemegang hak cipta.

\medskip
\noindent Seluruh kode, model, dan berkas konfigurasi yang dibahas tersedia pada \url{https://github.com/alfa-yohannis/advanced-software-engineering}.
\clearpage
